% NOTES:
% manually added german quotationmarks: \glqq \grqq // or use UTF-8 characters directly: „quotation marks“
% automatically added quotationmarks: \enquote{Text} // also supports nesting: \enquote{Test with \enquote{nested} quotation.}
% citation: \cite{Pohl2005} -> [1] /// \cite[S. 9ff]{Pohl2005} -> [2, S. 9ff]
%           (ISO-)Standards: http://itooktheredpill.dyndns.org/2010/citing-iso-standards-in-latex-using-bibtex/
% use of "et al.": http://blog.apastyle.org/apastyle/2011/02/et-al-when-and-how.html
%                  First occurence: Number of authors >= 6: use "et al.": "Sunderland et al. explained that.."
%                                   Number of authors <  6: name surnames of all authors
%                  Further occurence: Number of authors >= 3: use "et al.": "Sunderland et al. explained that.."
% references: Add label with \label{label1} and then refer to it with e.g. \hyperref[label1]{Label \ref*{label1} on page \pageref*{label1}}
%                                                                        or conveniently via \myRef{Abschnitt}{label}
% bibliography: keep information minimal (no URLs, no DOI, no ISBN)
% acronyms: \ac{KDE}  == K Desktop Environment (KDE) // on first occurence; on further use only acronym
%           \acs{KDE} == KDE // plural: \acsp{KDE} == KDEs
%           \acf{KDE} == K Desktop Environment (KDE) // \acfp{KDE} == K Desktop Environments (KDEs) 
%           \acl{KDE} == K Desktop Environment // plural: \aclp{KDE}
%       (!) \acp{KDE} == K Desktop Environments (KDEs) (on first occurrence) == KDEs (on further occurrences)
%           Definition: \acro{DD}{Data-Dictionary} // standard acronym definition
%                       \acrodefplural{DD}{Data-Dictionaries} // special plural definition (usually only an 's' is appended: \acp{KDE} == KDEs)
%                       Usage: \acl{DD} == Data-Dictionary // \aclp{DD} == Data-Dictionaries
% non-breaking space: used within abbreviations like z.\,B. to prevent both parts from being separated (e.g. due to line breaks)
%                     \, creates a narrow non-breaking space e.g. used for dates: 23.\,05.\,2042
%                     ~ creates a full non-breaking space e.g. used for full dates (23.\,Mai~2012)
%                        or terms that must not be separated (Prof.\,Dr.\,Bruno~Müller-Clostermann)
% wording: prefer the term "part" or "paragraph" to "chapter" unless you are writing a book or PhD thesis ;-) (German: "Abschnitt" statt "Kapitel" verwenden)
% LaTeX-Tabus: ftp://ftp.dante.de/tex-archive/info/german/l2tabu/l2tabu.pdf 
%
%
% PROBLEMS AND SOLUTIONS:
% Package hyperref Warning: Token not allowed in a PDFDocEncoded string... removing `math shift' on input line XXX.
%    Solution: provide alternative plain text for mathsymbols or special characters that will end up in links. 
%    E.g. change \section{The joy of $\pi$} to \section{The joy of \texorpdfstring{$\pi$}{pi}}
%    http://homepages.inf.ed.ac.uk/imurray2/compnotes/latex.html
% Biliography is not created properly
%    Problem: References are not found after building, log file and \printbibliography command show messages such as "empty bibliography".
%    Try to run biber manually on command line. Possible biber caching error (encountered on OS X 10.8.2):
%    "data source /var/folders/m6/bn7r45zx6cx55rr9g4qh6s6w0000gn/T/par-agoldst/<..>/inc/lib/Biber/LaTeX/recode_data.xml not found in ."
%    Solution: sudo rm -rf /var/folders/m6/bn7r45zx6cx55rr9g4qh6s6w0000gn/T/par-agoldst/
%    see http://humtex.wordpress.com/2011/11/29/biber-first-aid-for-data-source-not-found/
%
%

% ########################### DOCUMENTCLASS ###########################

\documentclass[paper=a4, % equivalent to a4paper
               headsepline, % line after headline
               fontsize=12pt,
               headings=small,
               %draft, % debug mode
               oneside, % one sided document: no alternating page number positions and margins (equivalent to twoside=false)
               ngerman, % new German spelling rules (important for hyphenation)
               captions=tableheading, % within table environment: put \caption{} text above table
               toc=bibliographynumbered, % adds bibliography to table of contents (numbered like chapters; unnumbered with toc=bibliography)
               1.5headlines]
               {scrreprt}

% ######################### PACKAGES + CONFIG #########################

\usepackage[T1]{fontenc} % German umlauts, special chars, etc.
\usepackage[utf8]{inputenc} % input encoding (preferably utf8 with utf8 encoding within your editor! ansinew, latin1, etc. are also possible but require you to use cumbersome workarounds to write German umlauts e.g. "u for writing an ü)
\usepackage[ngerman]{babel} % new German spelling rules (important for hyphenation)
\usepackage[german]{varioref} % for commands like \vref to create references like "...see figures 3.4 to 3.6 on pages 23–24..."


% special symbols
\usepackage{pifont}

% strings have to be entered after utf8 has been set as input encoding

% strings for the frontpage
\newcommand{\documentType}{Bachelorarbeit} % Art der Arbeit: Master-/Bachelor-/Seminararbeit/etc.
\newcommand{\documentHeading}{Natural Language Processing zur Verbesserung von Suchfunktionen in der Software-Entwicklung}
\newcommand{\documentLocation}{Essen} % also used as location for declaration of academic honesty / Ort in Eidesstattlicher Erklärung
\newcommand{\documentDate}{\today}%10.\,04.\,2014} % also used as date for declaration of academic honesty / Datum in Eidesstattlicher Erklärung

% PDF document metadata (the following strings must not contain line breaks, non-breaking spaces, etc.)
\newcommand{\documentPDFauthor}{Leon Zimmermann}
\newcommand{\documentPDFtitle}{\documentHeading}



\usepackage{acronym} % enables acronym environment ftp://ftp.tu-chemnitz.de/pub/tex/macros/latex/contrib/acronym/acronym.pdf
                     % http://de.wikibooks.org/wiki/LaTeX-W%C3%B6rterbuch:_Abk%C3%BCrzungsverzeichnis

\usepackage{hyphenat} % hyphenation
\hyphenation{} % special hyphenation for certain words (you can also simply add words that must not be hyphenated)

\usepackage{eurosym} % Euro currency symbol: use \EUR{10} or 10 \euro (\EUR{10} recommended because of correct space between number and symbol)
\DeclareUnicodeCharacter{20AC}{\euro} % UTF-8: enables use of Euro currency sign € (cf. http://www.theiling.de/eurosym.html)

\usepackage{amsmath, amsthm, amssymb} % AMS math symbols and theorem boxes; configure via \newtheoremstyle (cf. ftp://ftp.ams.org/pub/tex/doc/amscls/amsthdoc.pdf)
\usepackage{mathtools}

\usepackage[top=2.5cm, bottom=2cm, left=5cm, right=1.4cm, headsep=1cm, footskip=1cm]{geometry} % borders
%\renewcommand{\restoregeometry}{\newgeometry{top=2.5cm, bottom=2cm, left=5cm, right=1.4cm, headsep=1cm, footskip=1cm}}
    % examples from http://www.artofproblemsolving.com/Wiki/index.php/LaTeX:Layout
    % headsep  = Distance from bottom of header to the body of text on a page.
    % footskip = Distance from bottom of body to the bottom of the footer.
\usepackage{scrlayer-scrpage} % header/footer https://ctan.org/pkg/scrlayer-scrpage?lang=de

\usepackage{setspace} % enables line spacing via \singlespacing, \onehalfspacing, \doublespacing or \setstretch{<factor>} for custom spacing
\usepackage{verbatimbox}
\usepackage{rotating} % rotate objects e.g. with \rotatebox{90}{}; also used for tables in landscape format with sidewaystable environment
\usepackage{tabularx}
\usepackage{calc}
\usepackage{array}
\usepackage[font={bf, small, sf}]{caption} % customize the captions in floating environments
\usepackage{subcaption}
\usepackage[final]{pdfpages} % inclusion of external PDF documents (options e.g.: selected pages including a blank one: pages={3,{},8-11,15} )
\usepackage{multirow} % span column over multiple rows
\usepackage{booktabs} % table rules with various thicknesses (\toprule, \midrule, \bottomrule, etc.)
\usepackage{longtable, tabu} % tables that continue to the next page; http://mirror.unl.edu/ctan/macros/latex/contrib/tabu/tabu.pdf 
\usepackage{ltcaption} % caption for longtable
\usepackage{multicol} % add areas with multiple columns
\usepackage{enumitem} % prefix for enumerations
\setlistdepth{9} % enumitem: enable long nested lists
\setlist[itemize,1]{label=$\bullet$}
\setlist[itemize,2]{label=$\bullet$}
\setlist[itemize,3]{label=$\bullet$}
\setlist[itemize,4]{label=$\bullet$}
\setlist[itemize,5]{label=$\bullet$}
\setlist[itemize,6]{label=$\bullet$}
\setlist[itemize,7]{label=$\bullet$}
\setlist[itemize,8]{label=$\bullet$}
\setlist[itemize,9]{label=$\bullet$}
\renewlist{itemize}{itemize}{9}

%\bibliographystyle{alpha} % examples: http://amath.colorado.edu/documentation/LaTeX/reference/faq/bibstyles.html
%\bibliographystyle{dinat} % bibliographies following the german norm DIN 1505
\usepackage[bibstyle=numeric,
            citestyle=numeric,
            backend=biber,
            maxcitenames=2,
            maxbibnames=100,
            sortcites=true,
            doi=false,
            isbn=false,
            url=false]{biblatex} % bibliography package; replaces old natbib package; numeric = Reference in brackets [1].
                                        % remember to translate via "biber file.bcf" (or configure LaTeX editor)
                                        % very useful for printing parts of the bibliography within the document
                                        % see http://www.komascript.de/node/1220 for filter examples
                                        % http://tex.stackexchange.com/questions/32851/biblatex-combine-filters
                                        % http://biblatex.dominik-wassenhoven.de/download/DTK-2_2008-biblatex-Teil1.pdf
%\nocite{*} % display all entries in bibliography regardless of whether they have been cited or not
\DefineBibliographyStrings{ngerman}{andothers={et~al\adddot}} % "u.a." zu "et al."

\usepackage[german=quotes]{csquotes} % dynamic and configurable quotationmarks that \enquote{also support \enquote{nested} quotations.}


\usepackage{lmodern} % Latin Modern fonts
%\usepackage{picins} % insert pictures into paragraphs

\usepackage{float}
\usepackage{wrapfig} % wrap text around figures with wrapfigure environment (has to be loaded right after float package)
\restylefloat{figure} % float only at the place where you define it (without LaTeX moving it)
\floatstyle{plain} % standard: no box is drawn


\usepackage{color} % define own colors: \definecolor{MyColor}{rgb}{0.37,0.17,0.97}
\usepackage{xcolor}
\usepackage{colortbl} % color tables and cells
\usepackage{adjustbox} % http://mirror.unl.edu/ctan/macros/latex/contrib/adjustbox/adjustbox.pdf

\usepackage{framed} % framed or shaded regions that can break across pages

\usepackage{graphicx} % enhanced support for graphics

\usepackage{url} % e.g. \url{http://www.example.com}

\usepackage{fullwidth} % full width environment to display something beyond the text width http://www.ctan.org/tex-archive/macros/latex/contrib/fullwidth


% hyperref package has to be loaded after the float package
\usepackage[%pdftex=true, % only if you use pdftex
            hypertexnames=true,
            pdfstartview={FitH}, % fit the width of the page to the window
            colorlinks=true,
            linkcolor=black,
            citecolor=black,
            filecolor=black,
            urlcolor=black]{hyperref} % clickable urls in PDF: \href{URL}{text}
\hypersetup{% configure PDF metadata with \hypersetup command to pass UTF-8 characters (e.g. umlauts) properly
            pdftitle={\documentPDFtitle}, % title (PDF metadata)
            pdfauthor={\documentPDFauthor}, % author (PDF metadata)
}
\urlstyle{tt} % font style of url: rm = CM Roman; sf = CM Sans Serif; tt = CM Typewriter; same = selected font


% Pseudo Sourcecode Formatting with algorithmicx (http://ctan.org/pkg/algorithmicx)
%   Packages have to be loaded after float and hyperref
%   Alternative: algorithm2e which cannot split code on multiple pages but has nice vertical lines
%   Get vertical lines in algorithmicx as seen in algorithm2e: http://tex.stackexchange.com/questions/110431/ploblems-with-vertical-lines-in-algorithmicx
%   Chaining procedure calls: http://tex.stackexchange.com/questions/16046/how-to-nest-call-in-algorithmicx
\usepackage{algpseudocode} % part of algorithmicx
\usepackage[Algorithmus]{algorithm} % float environment algorithm, parameter is the prefix title in every algorithm box (e.g. "Algorithm 1")


% margin notes
\setlength{\marginparwidth}{1.2in} % width of margin notes
\reversemarginpar % margin notes (added via \marginpar{margin text}) are placed on the opposite side (left/inside instead of right/outside)


% load custom commands and strings
% suppresses page numbering in the header for one page (usage e.g.: \suppressPagenumber{\part{Stand der Forschung}})
\newcommand{\suppressPagenumber}[1]{{\ohead[]{} #1 \ohead[\rightline{\thepage}]{\rightline{\thepage}}}}


% own quotation style for a big single quotation within a block of text
\newcommand{\myQuotation}[2]{\begin{quote}\small{\textit{\enquote{#1}}} -- #2\end{quote}}



% side notes with marginpar
\newcommand{\myNote}[1]{\marginpar{\raggedleft \footnotesize \color{gray} \textit{#1}}}



% cite reference with full title ahead
\newcommand{\myCiteTitle}[1]{\citetitle{#1} (\citeauthor{#1}~\cite{#1})}




% Image/Figure
% Usage: \myImg{caption}{label}{width-factor}{path}
%   Example: \myImg{Beispiel eines Datenflussdiagramms}{fig:dfd}{1}{data-flow-diagram.pdf}
% http://www2.informatik.hu-berlin.de/~piefel/LaTeX-PS/Archive-2004/V08-grafikein.pdf
\newcommand{\myImg}[4]{
    \begin{figure}[h!]
        \centering
        \includegraphics[width=\textwidth * \real{#3}]{./img/#4}
        \caption{\label{#2}#1}
    \end{figure}    
}
\newcommand{\myImgCfg}[5]{
    \begin{figure}[#5]
        \centering
        \includegraphics[width=\textwidth * \real{#3}]{./img/#4}
        \caption{\label{#2}#1}
    \end{figure}    
}

% Shortcut for "..Abbildung 2 zeigt.." via "..\myImgRef{fig:screenshot} zeigt.."
\newcommand{\myImgRef}[1]{\hyperref[#1]{Abbildung~\ref*{#1}}}



% Shortcut for "..Abschnitt 2 beschreibt.." via "..\myRef{Abschnitt}{chap:myrefchapter}.."
\newcommand{\myRef}[2]{\hyperref[#2]{#1~\ref*{#2}}}


% Shortcut for "..auf Seite 7.." via "..\myPageRef{fig:myfigure}.."
\newcommand{\myPageRef}[1]{\hyperref[#1]{Seite~\pageref*{#1}}}



% Shortcuts for frequently used strings
\newcommand{\sczb}{z.\,B. } % zum Beispiel; added space because those phrases will never be used at the end of a sentence
\newcommand{\scdh}{d.\,h. } % das heißt
\newcommand{\scua}{u.\,a. } % unter anderem






% Mathematics
\newcommand{\powerset}[1]{\mathcal P~\left({#1}\right)}
%\newcommand{\powerset}[1]{\wp \left({#1}\right)}



% Font
% changes font family ptm = Times, phv = Helvetica,.. http://de.wikibooks.org/wiki/LaTeX-W%C3%B6rterbuch:_fontfamily
%\renewcommand{\familydefault}{phv}



% colored box
% see LaTeX - fcolorbox.png
% http://www.lessjunkmorefunk.de/de/node/10
%\definecolor{Gray}{gray}{0.9}
%\newcommand{\inhalt}[1]{\fcolorbox{black}{Gray}{\parbox{\textwidth}{#1}}}


\addbibresource{./referenzen.bib} % ./ for local folder is very important!


% ############################# DOCUMENT ##############################
\begin{document}
\shorthandoff{"} % deactivates shorthand "a for a umlaut (ä) (necessary if quotation marks are used e.g. in source code)
\setcounter{secnumdepth}{3} % numbering up to three level
%\sloppy % LaTeX will always try to create justified text (this may result in bigger gaps between words)

% ############################# FRONTPAGE #############################
\pagestyle{empty} % remove page number

\includegraphics[width=\textwidth]{./img/logo_unidue_paluno_deutsch.pdf}

\begin{sffamily}

\begin{center}

\vspace*{0.7cm}

\normalsize{Institut für Informatik und Wirtschaftsinformatik (ICB)\\%Institute for Computer Science and Business Information Systems (ICB)
Lehrstuhl für Software Engineering, insb. mobile Anwendungen\\
Prof. Dr. Volker Gruhn}

%\vspace{\baselineskip}

\vspace*{2.5cm}

\LARGE{\documentHeading}

\vspace*{1cm}

\large{\documentType}

\vspace*{2.5cm}

\normalsize{vorgelegt der Fakultät für Wirtschaftswissenschaften\\
der Universität Duisburg-Essen (Campus Essen) von}

\vspace*{\baselineskip}

\normalsize{
Leon Zimmermann \\
Laddringsweg 8 \\
45219 Essen \\
Matrikelnummer: 3080384
}

\vspace*{1.5cm}

\small{\documentLocation, den \documentDate}

\vspace*{1.5cm}

\normalsize{
\begin{tabular*}{\linewidth}{p{0.3\linewidth}p{0.7\linewidth}}
\textbf{Betreuung:}      & Wilhelm Koop, Sascha Feldmann\\
\textbf{Erstgutachter:}  & Prof.\,Dr.\,Volker~Gruhn\\
\textbf{Zweitgutachter:} & ???\\
\\
\textbf{Studiengang:}    & Angewandte~Informatik - Systems~Engineering (B.\,Sc.)\\
\textbf{Semester:}       & 10
\end{tabular*}
}

\end{center}

\end{sffamily}
\clearpage % calls \pagestyle{empty}
%\onehalfspacing % set one-and-a-half line spacing after frontpage was included
\setstretch{1.1} % a value of 1.25 refers to one-and-a-half spacing and 1.667 to double-spacing (for 10pt)
                 % (LaTeX calculates with a leading of 1.2 at a point size of 10pt: 1.2 * 1.25 = 1.5 line spacing)

% ############################# HEADLINES ######################
\renewcommand{\headfont}{\small} % small headline font size to fit even long headlines onto the page

\clearscrheadfoot
\ohead[\rightline{\thepage}]{\rightline{\thepage}} % headline with number of page on the right hand side
\automark{chapter}

\pagenumbering{roman} % roman numbering: I, II, III, IV, ...

\pagestyle{scrheadings} % defines whether headlines are printed and how they are styled

% ############################## ABSTRACT #############################

%- Was ist das generelle Thema der Arbeit?
%Beispiel: Kontext ist der Schnittpunkt Variabilitätsmodellierung / SPLE / SA.
%
%- Was ist das ganz konkrete Problem, das in der Arbeit behandelt wird? Was ist die Forschungslücke?
%Beispiel: DFDs haben keine Variabilitätserweiterung. Bisherige Lösungen behandeln das Thema nur abstrakt und können eine Konsistenzprüfung im Sinne des SPLE-Gedankens nicht konkret umsetzen.
%
%- Was ist die generelle Vorgehensweise (Methodik, z.B. Art der Studie: Experiment, Literatursuche, …)?
%Beispiel: SLR zum Aufspüren von Problemen und potentiellen Lösungen mit dem Ziel, diese zu einem eigenen Ansatz zu kombinieren.
%
%- Was ist das zentrale Resultat? Was folgt aus dem Resultat (Wert der Arbeit, z.B. für künftige Studien)?
%Beispiel: Ansatz zur SA im SPLE mit Toolunterstützung. Macht den Weg frei für weiterführende Ansätze in dem Rahmen (evtl. ein knackiges Beispiel nennen, was man tun kann).
%
%Dabei sollte die Wortwahl – im Gegensatz zu den von mir gegebenen Stichpunkten – natürlich durchgehend knackig und äußerst präzise sein. Zu jedem Punkt sollten jeweils ein bis zwei Sätze formuliert werden.
%Insgesamt sollte ein Abstract ca. 100-250 Wörter lang sein (grob eine eine halbe Seite bei dem üblichen Format, niemals mehr als eine Seite).
%
%Du solltest dabei daran denken, dass ein Abstract potentiellen Lesern zur Grobeinordnung Deiner Arbeit dient, z.B., um festzustellen, ob die Arbeit für ihre Forschung relevant sein könnte. Außerdem ist ein Abstract das Aushängeschild der Arbeit; er sollte bei am Problem interessierten Lesern (auch bei Gutachtern!) großes Interesse wecken.
%%
%
%
%
\section*{Abstract}
% TODO: Zusammenfassung auf Englisch

%
%
%
\section*{Zusammenfassung}
Softwareentwickler müssen sich bei ihrer Arbeit Informationen zusammensuchen, welche sie für die weitere Arbeit benötigen.
So muss ein Softwareentwickler bei der Implementierung eines Features die Intention des Features kennen.
Solche Informationen können in Wissensdatenbanken hinterlegt sein.
Um dort die gewünschten Informationen zu finden können Suchfunktionen verwendet werden.
Insbesondere, wenn dem Softwareentwickler im Vorfeld nicht klar ist, wo er die gewünschten Informationen in der Wissensdatenbank finden kann.
Aber nicht immer liefert diese Suchfunktion die gewünschten Informationen.\\

In dieser Arbeit wird eine semantische Suchfunktion entwickelt.
Diese soll eine Verbesserung zu bestehenden Suchfunktionen bieten, wie beispielsweise die Suchfunktion von Confluence.
Um festzustellen, ob die neue Suchfunktion \textit{besser} ist als die bestehende, werden Methoden und Kriterien zur Evaluierung von Suchfunktionen erörtert.
Anhand dieser Kriterien werden die neu implementierte Suchfunktion und die bestehende Suchfunktion von Confluence in einer Studie verglichen.
In der Studie werden beispielhafte Sucheingaben definiert, sowie die erwarteten Ergebnisse.
Die Definition der Sucheingaben erfolgt auf Basis von Anwendungsfällen.
Die Anwendungsfälle beschreiben Situationen, in denen es realistisch ist, dass ein Softwareentwickler die Suchfunktion einer Wissensdatenbank verwendet.
Anhand von Argumentationen werden entsprechend realistische Sucheingaben erstellt.\\

% TODO: Ergebnis der Studie darstellen

% ########################### TOC + ACRONYMS ##########################
\setcounter{chapter}{0}

\setcounter{tocdepth}{2} % table of contents depth: 3 means e.g. 4.3.2.1 will be listed, whereas 2 will lead to 4.3.2. as the deepest entry, etc.
{
\sloppy
\tableofcontents 
}

\listoffigures

\listoftables

\chapter*{Abkürzungsverzeichnis}

\begin{acronym}
    \acro{ANN}{Approximate Nearest Neighbor}
    \acro{HNSW}{Hierarchical Navigable Small Worlds}
    \acro{NLP}{Natural Language Processing}
    \acro{VSM}{Vector Space Model}
\end{acronym}
\clearpage % Clear page to get the arabic page counting on the next page

% ########################## LAYOUT + DESIGN ##########################

\clearscrheadfoot % cleans footer from previous configurations
\ihead[]{\leftline{\leftmark}} % current section and chapter heading in header on the left
\ohead[\rightline{\thepage}]{\rightline{\thepage}} % page number in header on the right
\automark{chapter} % for creating the current section/chapter in the heading

\pagenumbering{arabic} % arabic numbering: 1, 2, 3, 4, ...
\setcounter{page}{1}

\addtokomafont{chapter}{\LARGE}
\addtokomafont{section}{\Large}
\addtokomafont{subsection}{\large}
\addtokomafont{paragraph}{\itshape}

\renewcommand{\thepart}{}%\Alph{part}}
\renewcommand{\thesubsubsection}{\Roman{subsubsection}}

\renewcaptionname{ngerman}{\partname}{}

% ############################## CONTENT ##############################

\chapter{Einleitung}
Für die tägliche Arbeit benötigt ein Softwareentwickler Informationen, welche über den Code, an dem er arbeitet, hinausgehen.
Um an diese Informationen zu kommen kann der Softwareentwickler eine Person suchen, welche ihm die gewünschte Information geben kann.
Diese Person kann er im Projekt finden, oder über Websites, wie StackOverflow.
Darüber hinaus wird in vielen Projekten eine Wissensdatenbank angelegt, welche Informationen enthält, welche spezifisch für das Projekt sind.
Eine solche Wissensdatenbank ist Confluence.
Sie bietet eine Suchfunktion, welche es dem Softwareentwickler erleichtern soll, die gewünschten Informationen zu finden.
Beispiele für Informationen sind die Spezifikationen oder Dokumentationen eines Teiles der Software, mit welcher der Softwareentwickler gerade arbeitet.
Oder aber auch Best-Practices, Guides, oder Informationen darüber, wie die Software gestartet oder ausgeliefert wird.
Auch Informationen über den Projektplan sind für einen Softwareentwickler von Bedeutung.\\

Aber nicht immer finden Softwareentwickler die gewünschten Informationen mithilfe der Suchfunktion.
Ziel dieser Arbeit soll es sein, mithilfe von Wissen über Suchalgorithmen und Algorithmen aus dem Natural Language Processing zu zeigen, wie sich bestehende Suchfunktionen verbessern lassen.
Software, wie chatGPT\footnote{https://openai.com/blog/chatgpt} zeigt, dass Large Language Models eine valide Möglichkeit zur Information Extraction sind.
Es ist denkbar, dass Suchfunktionen für Softwareentwickler durch Large Language Models verbessert werden können.
Diese Arbeit erörtert, wie die Technologien von Large Language Models verwendet werden können, um eine semantsiche Suche zu implementieren.
Die implementierte semantische Suche wird mit der bestehenden Confluence-Suche verglichen, um herauszufinden, ob die semantische Suche eine Verbesserung zu der bestehenden Confluence-Suche darstellt.

\section{Vorgehensweise}
Die Gründe, warum eine gewünschte Information schwierig zu finden ist, sind vielzählig.
Manchmal kennt der Softwareentwickler nicht das genaue Wording, um die gewünschten Informationen zu finden.
Manchmal ist das Abstraktionslevel der gefundenen Informationen nicht das, welches sich der Softwareentwickler gewünscht hat.
Beispielsweise, wenn eine allgemeine Definition von Domänenobjekten gesucht wird, aber eine Spezifikation eines Anwendungsfalls gefunden wird, in welchem das Domänenobjekt lediglich erwähnt wird.
Im Rahmen dieser Arbeit wird eine neue Suchfunktion entwickelt, welche eine Verbesserung zu bestehenden Suchfunktionen von Wissensdatenbanken, wie Confluence, darstellen soll.
Nachdem die neue Suchfunktion entwickelt wurde, wird diese mit einer bestehenden Suchfunktion verglichen, um zu untersuchen, ob die neue Suchfunktion besser ist als eine bestehende Suchfunktion.
Dazu werden objektive Metriken zur Messung von Suchfunktionen dargestellt.
Die einzelnen Schritte für die Entwicklung und Evaluierung der neuen Suchfunktion werden im Folgenden dargestellt:
\begin{itemize}
   \item \textbf{Definition von Anwendungsfällen:}
         Es werden zuerst Anwendungsfälle definiert.
         Die Anwendungsfälle beschreiben die konkreten Situationen, in welchen ein Softwareentwickler eine Suche nutzen könnte.
         Das hilft später dabei Suchfunktionen miteinander zu vergleichen.
         Denn anhand der Anwendungsfälle können realistische Sucheingaben definiert werden.
         Diese Sucheingaben können an zwei verschiedene Suchfunktionen übergeben werden.
         Anschließend können die gefundenen Ergebnisse verglichen werden.
         Die genaue Vorgehensweise für diesen Vergleich wird in \myRef{Kapitel}{chap:vergleich-der-suchfunktionen} erklärt.
   \item \textbf{Technologien von Confluence:}
         Als nächstes wird die Technologie von Confluence dargestellt.
         Das hilft dabei später die Konzeption einer neuen Suchfunktion zu entwickeln und zu begründen, warum die neue Technologie besser sein soll.
         Im weiteren Verlauf der Arbeit wird erläutert, dass für den Vergleich der Suchfunktionen nicht die tatsächliche Confluence-Suchfunktion als Benchmark für die neue Suchfunktion herangezogen wird.
         Stattdessen wird die neue Suchfunktion so entwickelt, dass sie so konfiguriert werden kann, dass sie eine gute Heuristik für die Confluence-Suchfunktion darstellt.
         Dieses Kapitel erfüllt ebenfalls den Zweck eine vergleichbare Suchfunktion entwickeln zu können, welche anhand der objektiven Kriterien gemessen werden kann.
         Auch bei der Auswertung der Ergebnisse der Suchfunktion sind die Inhalte aus diesem Kapitel relevant.
   \item \textbf{Konzeption der neuen Suchfunktion:}
         Nachdem die Technologien von Confluence erörtert wurden, wird in diesem Kapitel die Theorie für die Implementierung einer Suchfunktion erläutert.
         Es wird zuerst erklärt, dass die neue Suchfunktion eine semantische Suchfunktion ist.
         Dann wird erläutert, was eine semantische Suche ist, und warum eine semantische Suche eine Verbesserung gegenüber den Technologien von Confluence darstellen soll.
         Anschließend wird erklärt, welche Technologien für die Implementierung einer semantischen Suche notwendig sind.
         Zuletzt wird die Implementierung der semantischen Suche mithilfe von der Vektordatenbank Weaviate\footnote{https://weaviate.io/} erklärt.
   \item \textbf{Vergleich der Suchfunktionen:}
         In diesem Kapitel werden zuerst Methoden für die Evaluierung von Suchfunktionen herausgearbeitet.
         Damit wird die Frage beantwortet, wann eine Suchfunktion \textit{gut} ist.
         Das hier erläuterte Wissen wird für die Durchführung der Studie benötigt.
         Denn um festzustellen, ob die Implementierung eine Verbesserung darstellt, muss eine Studie durchgeführt werden.
         Aufgrund des Scopes der Arbeit wird nur eine rudimentäre Studie durchgeführt.
         Die Studie vergleicht die Precision-Werte der Confluence-Suchfunktion und der neuen Suchfunktion mithilfe eines generierten Datensatzes von Sucheingaben und erwarteten Ergebnissen.
         Grundlage für die Generierung dieses Datensatzes sind die Anwendungsfälle, welche zuvor definiert wurden.
         Die Ergebnisse werden dargestellt und diskutiert.
         Es wird gezeigt, dass eine semantische Suche für die generierten Daten und den spezifischen Kontext der Datengrundlage, einen schlechten Precision-Wert besitzt. 
         Dann werden dieses Ergebnis interpretiert und es wird der Schluss gezogen, dass Sentence Transformer besser geeignet sind für Open-Domain Datensätze, im Gegensatz zu Closed-Domain Datensätzen.
         Zuletzt wird der Studienaufbau diskutiert.
    \item \textbf{Zusammenfassung und Ausblick}
         In diesem Kapitel werden die einzelnen Kapitel der Arbeit nochmal zusammenfassend dargestellt.
         Außerdem wird das Ergebnis der Arbeit zusammenfassend dargestellt.
         In \myRef{Kapitel}{chap:ausblick} werden Ansätze erörtert, welche die Suchergebnisse von semantischen Suchen in spezifischen Domänen verbessern sollen.
\end{itemize}

\section{Verwandte Arbeiten}
Es gibt einige Arbeiten, welche die gleichen oder sehr ähnliche Probleme adressieren.
Die verschiedenen Herangehensweisen werden in diesem Kapitel erläutert.
Zuerst wird der Begriff Traceability behandelt, und wie diese genutzt wird, um Suchfunktionen zu verbessern.
Die Herangehensweisen zur Herstellung von Traceability umfassen die Verwendung von Word Embeddings oder auch ein Task-Based Ansatz.
Es folgt eine Übersicht über Arbeiten, welche sich mit Question Answering beschäftigen.
Dann folgt eine Übersicht über Arbeiten, welche sich mit der Evaluierung von Suchfunktionen auseinandersetzen.

% Traceability
Traceability bedeutet, dass sich von einer Stelle im Code, auf die entsprechenden Stellen in anderen Artefakten zurückschließen lässt.
Ein Anwendungsfall für eine solche Traceability-Funktionalität ist \textit{Feature Location}, also das Finden der Spezifikation eines Features, wenn nur der Code vorhanden ist.
Analog dazu ist \textit{Bug Localization} ein Anwendungsfall zum Auffinden von Stellen im Code, welche mit einem Bug zusammenhängen.
Zur Herstellung von Traceability zwischen Code und Dokumentation gibt es verschiedene Ansätze.
Haiduc et. al. schlagen ein System zur Verbesserung von Sucheingaben vorgeschlagen (\cite{Haiduc_Bavota_Marcus_Oliveto_DeLucia_Menzies_2013}).
Das System verwendet Query Reformulations, um die Traceability herzustellen.
Query Reformulation bedeutet, dass das System den Softwareentwickler bei der Eingabe einer Suchanfrage zur Suche nach den passenden Artefakten unterstützt.
Dazu gibt der Softwareentwickler zunächst eine Suchanfrage ein, und markiert diejenigen Ergebnisse, welche am relevantesten für ihn sind.
Auf Grundlage der gewählten Ergebnisse und mithilfe eines Machine Learning Algorithmus werden nun Vorschläge für eine verbesserte Suchanfrage gemacht.
Dabei gibt es verschiedene Strategien.
Wenn der Softwareentwickler zu Beginn eine sehr lange Suchanfrage eingegeben hat, dann kann das System eine Reduktion der Suchanfrage vorschlagen.
Hat der Softwareentwickler dagegen lediglich einen Suchbegriff angegeben, so kann das System eine Erweiterung der Suchbegriffe vorschlagen.
Dazu greift das System auf Synonyme des eingegebenen Suchbegriffes zurück.\\

Ye et. al. beschreiben, wie Word Embeddings dazu verwendet werden können, um Traceability zwischen Code und anderen Softwareentwicklungs-Artefakten herzustellen.
Word Embeddings sind eine Datenstruktur, welche einem Wort einen Vektor in einem n-dimensionalen Raum zuweist.
Anhand dieses Vektors kann die Ähnlichkeit zwischen Wörtern beschrieben werden.
Ähnliche Wörter haben eine geringe Distanz im n-dimensionalen Raum.
Unähnliche Wörter haben eine hohe Distanz.
Der Algorithmus, welcher die Ähnlichkeit der Wörter bestimmt, macht Gebrauch von der Distributional Hypothesis (\cite{Harris_1954}).
Dieser besagt, dass Wörter, welche im gleichen Kontext verwendet werden, eine ähnliche Semantik besitzen.
Hermit wird also die Ähnlichkeit der Wörter bestimmt.
Dieses Verfahren wird nun sowohl auf den Code angewendet als auch auf die Softwareentwicklungs-Artefakte (\cite{Ye_Shen_Ma_Bunescu_Liu_2016}).\\

Antoniol et. al. verwenden einen ähnlichen Ansatz, wie Ye et. al.
Auch hier werden Word Embeddings verwendet um Softwareentwicklungs-Artefakte gegen den Code zu matchen.
Hier durchlaufen die Artefakte und der Code zwei verschiedene Pipelines.
Die Wörter der Artefakte in natürlicher Sprache werden in lowercase umgewandelt.
Anschließend werden Stoppwörter entfernt.
Zuletzt werden Flexionen entfernt.
Aus dem Code werden zunächst Identifier extrahiert.
Identifier, welche mehrere Wörter unter Verwendung von CamelCase oder snake\_case beinhalten, werden in die einzelnen Wörter aufgeteilt.
Anschließend werden die Identifier auf die gleiche Art und Weise normalisiert, wie die Wörter der Softwareentwicklungs-Artefakte.
Dann erfolgt sowohl für die Identifier als auch für die Wörter aus den Artefakten die Indizierung, also die Umwandlung in Word Embeddings (\cite{Antoniol_Canfora_Casazza_DeLucia_2000}).\\

Treude et. al. entwickeln eine Oberfläche, welche die Suche von \textit{Tasks} ermöglicht.
Damit soll ebenfalls Traceability zwischen Code und Dokumentation hergestellt werden.
Dabei ist unter Task eine Operation im Code zu verstehen.
Sie beschreiben einen Task als Verben, welche mit einem direkten Objekt oder einer Präposition in Verbindung stehen.
Die Autoren nennen die Phrasen \textit{get iterator} und \textit{get iterator for collection} als Beispiele.
Die Software analysiert nun die gesamte Dokumentation und extrahiert Tasks.
Die Tasks werden in einen Index geschrieben, sodass der Softwareentwickler nach ihnen suchen kann (\cite{Treude_Sicard_Klocke_Robillard_2015}).\\

% Question Answering
Neben der Traceability zwischen Code und Dokumentation ist das Question Answering ein Ansatz zur Verbesserung von Suchfunktionen.
Suchfunktionen sind Document Retrieval Systeme.
Sie liefern Dokumente, welche zu der Sucheingabe des Nutzers passen.
Document Retrieval Systeme sind eine Unterkategorie von Information Retrieval Systemen.
Information Retrieval Systeme liefern auf Anfrage Informationen an den Nutzer.
Im Fall einer Suchfunktion werden Dokumente geliefert, welche diese Information beinhalten.
Mithilfe der Dokumente ist der Ort, an dem sich die, vom Nutzer gewünschte, Information befindet eingegrenzt.
Nichtsdestotrotz muss der Nutzer aus dieser Eingrenzung die gewünschte Information manuell extrahieren.
Ganz im Gegensatz zu Question Answering Systemen.
Zhang et. al. untersuchen verschiedene Herangehensweisen zur Implementierung von Open-Domain Question Answering Systemen (\cite{Zhang_Chen_Xu_Cao_Chen_Cohn_Fang_2023}).
Open-Domain Question Answering Systeme beantworten allgemeine Fragen eines Nutzers, z.B. basierend auf Informationen von Wikipedia.
Closed-Domain Question Answering Systeme beantworten dagegen Fragen im Kontext einer spezifischen Domäne, z.B. basierend auf unternehmensinternen Informationen.\\
% TODO: closed-domain sentence transformer transformer zu verwandten arbeiten hinzufügen (fine-tuned transformers)

% TODO: Informationen lassen sich auch anreichern: NER etc. (NLP Techniken)
% TODO: Informationen lassen sich aufbereiten: RAG

% Messung von Performance
% TODO: Trec und Benchmarks hinzufügen
Für die Messung der Performance einer Suchfunktion gibt es verschiedene Metriken.
Metriken können systemspezifisch sein, nutzerspezifisch.
Systemspezifische Metriken messen die Performance anhand objektiver Kriterien, welche an dem System gemessen werden können.
nutzerspezifische Metriken messen die Performance dagegen anhand subjektiver Kriterien.
Die Messung erfordert Probanden, welche die Suchfunktion verwenden.
Die Probanden legen die Performance der Suchfunktion anhand der zu betrachtenen subjektiven Kriterien fest.
Clarke und Willet messen die Qualität von Suchfunktionen des World Wide Webs anhand des Recalls (\cite{Clarke_Willett_1997}).
Um dies zu ermöglichen wird ein Datensatz generiert, welcher Sucheingaben beinhaltet, sowie alle relevanten Dokumente für eine Sucheingabe.
Bar-Ilan verwendet die gleichen Metriken (\cite{Bar-Ilan_2002}).
Hier wird darüber hinaus auf die Problematik der Messung des Recalls eingegangen.
Es wird erläutert, dass zur Messung des Recalls a-priori bestimmt werden muss, welche Dokumente als relevant für eine gegebene Sucheingabe erachtet werden sollten.
Gordon und Pathak behaupten, dass die Bestimmung der Relevanz lediglich dem Nutzer mit dem Bedürfnis nach der Information zu überlassen ist (\cite{Gordon_Pathak_1999}).
Voorhees und Harman behaupten, dass die Bestimmung der Relevanz durch ein Experten-Panel durchgeführt werden sollte (\cite{Voorhees_Harman_2001}).
Sirotkin betrachtet verschiedene Ansätze zur Messung der Performance von Suchfunktionen (\cite{Sirotkin_2012}).
Neben den bereits genannten Metriken von Precision und Recall werden andere Metriken zur Messung der Performance einer Suchfunktion genannt, wie Mean Reciprocal Rank und Maximal Marginal Relevance.\\

\chapter{Auswahl von Use-Cases}

\section{Wahl des Abstraktionslevels}
Oft ist im Voraus nicht klar, was man genau sucht. Man gibt einen Suchbegriff ein und möchte "sinnvolle" Informationen zu diesem Suchbegriff bekommen. Sinnvoll könnte eine Definition des Begriffs sein, oder ein Überblick darüber, in welchen Bereichen und Kontexten dieser Begriff vorkommt.

TODO

\section{Semantik der Suche verstehen}
TODO

\section{Aktualität und Relevanz des Dokuments}
Manchmal werden Blog-Einträge von unbekannten Benutzern angezeigt, welche sehr kurz sind. Die Überschrift ist zwar passend, aber der Inhalt des Textes nicht. Und der Text ist nur sehr kurz. Es scheint offensichtlich, dass für allgemeine Suchanfragen längere Texte als Ergebnis sinnvoller sind, da eine allgemeine Suchanfrage der Versuch ist, sich mit einem Thema vertraut zu machen. In diesem Fall sollten viele Informationen priorisiert werden, gegenüber kurzen Texten, welche die Lösung eines sehr spezifischen Problems erläutern.

Manchmal werden alte Seiten gefunden oder unfertige Seite. Es sollte erkannt werden, wenn eine Seite veraltet ist (ist markiert), oder wenn sie noch nicht fertig ist (durch Markierung oder Länge/Struktur des Textes)

TODO
\chapter{Evaluationsmethoden und -Kriterien}

\section{Features}

\subsection*{Relevanz}
Eine Suchfunktion muss nicht nur Ergebnisse für eine Such-Query finden, sondern diese auch darstellen. Bei der Darstellung der Ergebnisse ist die Reihenfolge besonders wichtig. Manchmal hat der Nutzer keine genaue Vorstellung davon, was genau er sucht. Er möchte Informationen zu einem Thema finden, aber ohne dabei bestimmte Seiten oder Inhalte im Hinterkopf zu haben. Nichtsdestotrotz möchte der Nutzer mit minimalem Aufwand so viele wichtige Informationen, wie möglich erhalten. Die Sortierung von Ergebnissen der Suche soll genau das ermöglichen.
Nun ist zu klären, wie die Bestimmung der Relevanz von Suchergebnissen überhaupt funktioniert. Es muss bestimmt werden, anhand welcher Kriterien die Suchfunktion die Relevanz bestimmt. Für diese Kriterien gibt es wiederum unterschiedliche Quellen. Eine Quelle für Kriterien der Relevanz ist das Nutzerverhalten. Seiten, die besonders häufig besucht werden und auf denen lange verweilt wird, oder welche sehr häufig bearbeitet werden, sind gute Kandidaten für eine hohe Relevanz. Seiten, welche bei Nutzern praktisch in Vergessenheit geraten sind, scheinen nicht so wichtig zu sein, sodass die Relevanz für die Suche niedrig ist.   

\subsection*{Clustering}
Thematisch ähnliche Dokumente tendieren dazu relevant für die gleiche Suche zu sein. Aus diesem Grund ist es oftmals sinnvoll ähnliche Dokumente zu clustern, also zu gruppieren, und sie dem Nutzer als ein solches Cluster darzustellen. Alternativ kann auch ein Suchergebnis dargestellt werden und ähnliche Suchergebnisse werden angezeigt als „Ähnliche Suchergebnisse“.

\subsection*{Autovervollständigung}
Autovervollständigung bedeutet, dass während der Eingabe von Suchbegriffen bereits Ergänzungen für die gesamte Eingabe vorgeschlagen werden. Das ist vor allem ein hilfreiches Feature in Fällen, wo der Nutzer nur eine grobe Ahnung davon hat, was er eigentlich sucht. Mithilfe von Vorschlägen durch die Suchmaschine selbst kann der Nutzer seine Eingabe spezifischer fassen als er es ohne Vorschläge schaffen würde, sodass die Suchergebnisse besser zu dem passen, was er eigentlich sucht.

\subsection*{Best-Prefix-Highlighting}
Das ähnlichste Ergebnis zur Eingabe des Nutzers soll gehighlightet werden.

\subsection*{Synonyme}
Synonyme bei Suche mitbeachten.

\section{Qualitätskriterien}

\subsection*{Interaktivität}
Klassische Suchfunktionen sind so konzipiert, dass der Nutzer einen Suchbegriff eingibt, anschließend auf einen Submit-Button klickt und erst dann Suchergebnisse angezeigt bekommt. Der Nachteil dieser Vorgehensweise ist, dass der Nutzer erst spät Feedback darüber bekommt, ob er mit den eingegebenen Suchbegriffen auch Ergebnisse bekommt, die für ihn gerade sinnvoll sind. Das Qualitätskriterium Interaktivität gibt an, wie kurz die Zeit zwischen dem Eingeben von Suchbegriffen und der Anzeige von Ergebnissen ist. Eine Suchfunktion muss nicht unbedingt einen Submit-Button verwenden. Die Suche kann auch mit jedem ausgeschriebenen Keyword die Suche starten und Ergebnisse dafür anzeigen, oder auch nach jedem eingegebenen Buchstaben. Mit einer kürzeren Zeit zwischen Eingabe und Ergebnis, also einer höheren Interaktivität, weiß der Nutzer schneller, ob seine Eingaben sinnvoll sind, und bekommt damit schneller gute Suchergebnisse.

\subsection*{Geschwindigkeit}
Die Geschwindigkeit einer Eingabe ist die Zeit zwischen dem Beginn der Suche und der Anzeige der ersten Antwort. Die Geschwindigkeit hat nichts mit der Interaktivität zu tun. 

\subsection*{Darstellung}
Wie gut kann man anhand der Zusammenfassung der Ergebnisse den Inhalt bestimmen.

\subsection*{Benutzerfreundlichkeit}
Das System findet auch Ergebnisse für ähnliche Wörter. So sind bspw. Tippfehler kein Problem. Fuzzy-Search.

\section{Precision, Recall und F-Maß}
Zur Evaluation werden die Werte Precision und Recall gemessen. Der Precision-Wert ist das Verhältnis zwischen allen gefundenen Dokumenten und den gefundenen Dokumenten, die tatsächlich relevant sind. Der Recall-Wert gibt an, wie viele von den tatsächlich relevanten Dokumenten auch gefunden wurden. Es ist schwierig beide Werte zu optimieren, da der Precision-Wert versucht die Anzahl der gefundenen Dokumente einzugrenzen und der Recall-Wert versucht die Anzahl der gefundenen Dokumente zu erweitern. Das F-Maß fasst beide Werte zu einem neuen Wert zusammen.

\section{Befragungen}
TODO
\chapter{Such- und NLP-Algorithmen}

\section{Suchalgorithmen}

\subsection{Keyword Search}
TODO

\subsection{Phrase Search}
Eine Phrase Search ist die Suche nach Textausschnitten in Dokumenten.
Es werden also nicht mehrere Keywords einzeln betrachtet, sondern die Eingabe in das Suchfeld als Ganzen.
Anschließend wird in Dokumenten nach genau diesem Textausschnitt gesucht.

\subsection{Boolean Search}
TODO

\subsection{Field Search}
Eine Field Search sucht Dokumente anhand von Attributen.
Der Nutzer kann diese Attribute auswählen.
Wenn der Nutzer beispielsweise ein Dokument sucht, welches am 01.01.2005 erstellt wurde, dann kann die Eingabe der Suche so aussehen: "erstelldatum: 01.01.2005".
Es können beliebig viele Attribute verwendet werden, um die Suche einzugrenzen.\\

Neben der Verwendung der Attribute für die Suche selbst, können die Attribute komplementär zu einer anderen Art von Suche verwendet werden.
So kann eine Suchfunktion Buttons bereitstellen, über welche Filter festgelegt werden.
Nun kann eine Keyword Search durchgeführt werden, aber die gefundenen Dokumente werden mithilfe der Filter weiter eingeschränkt.

\subsection{Structured Search}
Eine klassische Suchfunktion verwendet Keywords, um relevante Dokumente für eine Sucheingabe zu ermitteln.
Der Vorteil dieser Art von Suche ist, dass die technische Struktur, in der die Daten vorliegen und gespeichert sind, nicht bekannt sein müssen.
Der Nachteil ist auf der anderen Seite, dass die Suchergebnisse unpräzise sein können.
Wenn beispielsweise der Suchbegriff Kamera eingegeben wird, dann werden Kameras als Ergebnisse zurückgegeben.
Soll die Kamera nun bestimmte Eigenschaften besitzen, dann müssen diese Eigenschaften ebenfalls als Keywords angegeben werden.
Nun wird aber nicht die Suche auf Ergebnisse eingegrenzt, bei denen eine Kamera diese bestimmten Eigenschaften besitzt.
Stattdessen werden Suchergebnisse angezeigt, bei denen einige dieser Keywords vorkommen.
Demgegenüber stehen Datenbankabfragen, beispielsweise mithilfe von SQL.
Bei einer Datenbankabfrage können Objekte abgefragt werden, dessen Eigenschaften ganz bestimmte Werte haben.
Die Ergebnisse, die eine solche Abfrage zurückgibt, sind dabei vollkommen genau.
Es werden keine Objekte zurückgegeben, welche diese Kriterien nicht erfüllen.
Voraussetzung für eine solche Suchabfrage ist allerdings, dass die Struktur der Datenbank a priori bekannt ist.
Dem Nutzer muss der Name der Datenbank, der relevanten Tabellen, sowie der relevanten Properties bekannt sein, damit er das passende SQL für die Datenbankabfrage schreiben kann.
Strukturierte Suchen sollen die Vorteile beider Vorgehensweisen kombinieren.
Die Struktur der Daten soll a priori nicht bekannt sein müssen, aber trotzdem sollen die Suchergebnisse vollkommen präzise sein.


\subsection{Lexical Search}
TODO

\subsection{Semantic Search}
Eine Semantische Suche arbeitet nicht anhand von Keywords, sondern anhand von Bedeutungen von Wörtern.
Sie versteht, dass einige Wörter sehr ähnlich sind, so wie "rollout" und "deployment", und dass diese Wörter oft im gleichen Kontext verwendet werden.
Um zu verstehen, welche Wörter kontextuell zusammengehören, werden die Wörter in einem n-dimensionalen Raum positioniert.
Wörter, die sich sehr ähnlich sind, also im gleichen Kontext verwendet werden, haben in diesem n-dimensionalen Raum eine geringe Distanz.
Wörter, die sich eher unähnlich sind, wie "rollout" und "API", haben eine größere Distanz.\\

Um eine semantische Suche zu implementieren, werden die Technologien von Transformern und Vektordatenbanken verwendet.
Ein Transformer bekommt als Input eine große Menge an Text und mappt die einzelnen Wörter auf einen Vektor einer beliebigen Länge.
Der Vektor, der am Ende herauskommt, beschreibt die Position des Wortes in dem n-dimensionalen Raum.
Der Vektor beschreibt gewissermaßen, wie stark ein Wort in eine abstrakte Kategorie einzuordnen ist.
Jeder Wert im Vektor entspricht einer Kategorie.
Mithilfe der Vektoren können verschiedene Wörter hinsichtlich ihrer Ähnlichkeit analysiert werden.
Ähnliche Wörter habe eine große räumliche Nähe, während zwei Wörter, die in vollkommen unterschiedlichen Kontexten verwende werden eine sehr große Distanz im Raum besitzen.
Nehmen wir für ein Beispiel einen dreidimensionalen Raum an.
Die X-Achse ist beschriftet mit dem Wort „Tier“, die Y-Achse ist beschriftet mit dem Wort „Computer“ und die Z-Achse ist beschriftet mit dem Wort „Mensch“.
Nun geben wir einem Transformer das Wort „Katze“, und der Transformer berechnet einen dreidimensionalen Vektor, welcher das Wort „Katze“ im Raum positioniert.
Weil eine Katze ein Tier ist, ist der X-Wert des Vektors eins.
 Der Wert eins bedeutet, dass das Wort vollständig zu dieser Kategorie gehört.
Da eine Katze überhaupt nichts mit einem Computer zu tun hat, ist der Y-Wert des Vektors 0.
Nun ist eine Katze zwar kein Mensch, aber eine Katze ist ein Haustier von Menschen.
Es ist denkbar, dass die Wörter Katze und Mensch oft im gleichen Kontext verwendet werden, sodass der Wert bei 0,3 liegen könnte.
Damit der Transformer einen Vektor berechnen kann, braucht er eine Menge Daten.
Diese Daten erhält er aus vielen Texten.
Werden zwei Wörter oft im gleichen Text genannt oder kommen zwei Wörter in vielen Texten sehr nahe beieinander vor, dann geht der Transformer davon aus, dass die beiden Wörter ähnlich sind, und berechnet ähnliche Vektoren.
Zuvor müssen die Texte allerdings bereitgestellt werden.
Dazu kann beispielsweise das Internet gecrawlt werden.
Die Ergebnisse des Transformers werden in einer Vektordatenbank gespeichert.
Eine Vektordatenbank ist eine Datenbank, welche Vector Embeddings, also ein Objekt als Key und dessen Vektor als Value speichert.
Bei dem Objekt kann es sich um Wörter handeln, dann wird auch von Word Embeddings gesprochen.
Es können aber auch Daten andere Daten, wie Bilder, Videos oder Audio gespeichert werden.
Der Zweck von Vektordatenbanken ist es, Daten nicht einfach linear zu speichern, sondern in einem Raum.
Die Distanz zwischen zwei Einträgen in diesem Raum beschreibt dessen Ähnlichkeit.
Genau diese Informationen machen sich semantische Suchen zu Nutze.

\section{NLP-Algorithmen}

TODO
\chapter{Vergleich der Suchfunktionen}
\label{chap:vergleich-der-suchfunktionen}
%TODO: Kapitel erweitern
%TODO: Testdaten äquivalente Beispiele darstellen
Zum Vergleich der beiden Suchfunktionen wird im Folgenden zuerst auf Evaluationsmethoden und -Kriterien eingegangen.
Diese sind die Grundlage für den Vergleich von zwei Suchfunktionen.
Anschließend wird der Aufbau der Studie dargestellt, sowie der zugrundeliegende Datensatz für die Durchführung der Studie.
Nachdem der Aufbau der Studie erklärt wurde, werden die Daten der Studie ausgewertet und dargestellt.
Das Ergebnis wird diskutiert.
Zuletzt wird auf die Validität der Daten eingegangen und es wird der Versuchsaufbau diskutiert.

\section{Evaluationsmethoden und -Kriterien}
\label{chap:evaluationsmethoden}
Um zu verifizieren, dass die Implementierung der neuen Suchfunktion ihre gewünschte Wirkung erzielt, müssen zuerst Methoden und Kriterien zur Messung herangezogen werden.
Dazu soll ein Experiment durchgeführt werden.
Bei einem Experiment wird zunächst eine Hypothese aufgestellt.
Diese wird anschließend durch das Experiment untersucht.
Ein solches Experiment wird für die Evaluation der Suchfunktionen erstellt.
Diese stützt sich auf den Anwendungsfällen, welche in \myRef{Kapitel}{chap:anwendungsfaelle} dargestellt wurden.
Dazu werden für die Anwendungsfälle Sucheingaben erstellt.
Für jede Sucheingabe wird definiert, welche Dokumente als Ergebnis erwartet werden.
Die Hypothese: Mit der implementierten Suchfunktion werden die erwarteten Dokumente \textit{besser} gefunden als mit der bisherigen Suchfunktion.
Diese Hypothese muss nun quantifizierbar und messbar gemacht werden.
Die subjektive Wahrnehmung reicht nicht für eine wissenschaftliche Arbeit aus.\\

Die Formulierung der Hypothese ist für sich genommen zu unspezifisch.
Es muss geklärt werden, wann eine Information \textit{besser} zu finden ist.
Im Folgenden werden Precision, Recall und F-Maß herangezogen, um Suchfunktionen anhand dieser Eigenschaften vergleichbar zu machen.
Mithilfe dieser Messwerte kann eine objektive Entscheidung darüber getroffen werden, welche Suchfunktion \textit{besser} funktioniert.

\subsection{Precision, Recall und F-Maß}
Zur Evaluation der Suchfunktionen können die statistischen Messwerte Precision und Recall verwendet werden.
Die Messwerte beschreiben, inwieweit eine Hypothese zutrifft.
Wenn also die Hypothese ist, dass ein bestimmtes Dokument gefunden wird, dann ist der Precision-Wert das Verhältnis zwischen allen gefundenen Dokumenten und den gefundenen Dokumenten, die tatsächlich relevant sind.
Der Wert lässt sich im Kontext der Suche nach Dokumenten wie folgt definieren (\cite{Sirotkin_2012}):\\

\(Precision=P=\frac{gefundene relevante Dokumente}{gesamte Anzahl gefundener Dokumente} \)\\

Der Recall-Wert gibt wiederum an, wie viele von den tatsächlich relevanten Dokumenten auch gefunden wurden.
Er lässt sich in diesem Kontext wie folgt definieren (\cite{Sirotkin_2012}):\\

\(Recall=R=\frac{gefundene relevante Dokumente}{gesamte Anzahl relevanter Dokumente}\)\\

Es ist schwierig beide Werte zu optimieren, da der Precision-Wert versucht die Anzahl der gefundenen Dokumente einzugrenzen und der Recall-Wert versucht die Anzahl der gefundenen Dokumente zu erweitern.
Das F1-Maß fasst beide Werte zu einem neuen Wert zusammen (\cite{Sirotkin_2012}):\\

\(F_1=2\frac{PR}{P+R}\)\\

Neben der Verwendung des F1-Maß müssen sich Gedanken darüber gemachen werden, welcher der beiden Messwerte wichtiger ist.
In diesem Fall ist es sinnvoll eher die Precision zu optimieren.
Denn, wenn ein Dokument gesucht wird, aber überhaupt nicht gefunden werden kann, dann erfüllt die Suchfunktion nicht ihren Zweck.
Wenn die Suchfunktion irrelevante Dokumente darstellt, kann sie trotzdem ihren Zweck erfüllen, solange die relevantesten Dokumente zuerst in der Liste der Ergebnisse dargestellt wird.
Dieser Faktor gilt auch bei der Implementierung zu berücksichtigen.

\section{Aufbau der Studie}
\label{chap:aufbau-der-studie}
Im Rahmen der Bachelorarbeit wird eine Studie durchgeführt, welche die neue Suchfunktion mit der bestehenden Confluence-Suche vergleicht.
Da der Zeitrahmen begrenzt ist, in welchem die neue Suchfunktion entwickelt wird, wird nicht die tatsächliche Suchfunktion von Confluence als Benchmark für die neue Suchfunktion herangezogen.
Das ist notwendig um einen gute Vergleichbarkeit zwischen den Suchfunktionen zu gewährleisten.
Die Produktreife von Confluence könnte die Performance dessen Suchfunktion positiv beeinflussen.
Da die neue Suchfunktion in nur kurzer Zeit entwickelt wurde, kann der Unterschied Produktreife ein Störfaktor bei dem Vergleich der Suchfunktionen sein.
Um dies zu vermeiden wird die Confluence-Suche repräsentiert durch eine BM25 Suche, welche ebenfalls mit Weaviate implementiert wurde. 
Es wurde bereits erwähnt, dass die Suche von Confluence auf Apache Lucene basiert, und dass dieses ein VSM auf Basis von TF-IDF verwendet.
Sparck et al. (\citeyear{Sparck_Jones_Walker_Robertson_2000}) beschreiben, dass eine große Ähnlichkeit bei der Performance zwischen TF-IDF und BM25 besteht.
Auf diese Weise wird die Vergleichbarkeit zwischen den beiden Suchfunktionen also verbessert.
Die BM25-Suche in Weaviate ist eine von drei Konfigurationen der neuen Suchfunktionen.
Sie ist die Benchmark für die neue Suchfunktion und repräsentiert die Confluence-Suche.\\

Eine weitere Konfiguration verwendet eine Mischung aus einer BM25 Suche, und einer semantischen Suche.
Die beiden Suchalgorithmen werden zu gleichen Teilen verwendet.
Zuletzt verwendet eine andere Konfiguration lediglich die semantische Suche auf Grundlage von Sentence Similarity.
Die Eigenschaft, welche untersucht werden soll ist, ob eine semantische Suche für den gegebenen Datensatz, und im Kontext einer Wissensdatenbank in der Softwareentwicklung, ebenfalls bessere Suchergebnisse liefert, als eine Suche auf Basis von BM25.
Durch die Verwendung einer einheitlichen Implementierung wird sichergestellt, dass lediglich die Unterschiede der Suchalgorithmen untersucht werden.
Es wird damit verhindert, dass die Confluence-Suche besser abschneidet, weil sie ausgereifter ist.
Es wird ebenfalls sichergestellt, dass die Datensätze der Suchfunktionen identisch sind.\\

Für die Studie wurde ein Datensatz generiert, welcher Dokumente beinhaltet, welche durch die Suchfunktion gefunden werden sollen.
Der Datensatz wurde aus einem realen Softwareentwicklungs-Projekt generiert, indem ein Teil der tatsächlichen Confluence-Seiten exportiert wurde.
Die Software soll im Gesundheitswesen eingesetzt werden.
Das bedeutet, dass die Wissensdatenbank Fachbegriffe aus dem Gesundheitswesen beinhaltet.
Darüber hinaus beinhaltet die Wissensdatenbank Definitionen, welche sich spezifisch auf das Softwareprojekt beziehen.
Daher handelt es sich bei dem Datensatz um eine Closed-Domain, also ein Datensatz, welcher spezifisch aus einer bestimmten Domäne generiert wurde.
Das steht im Gegensatz zu einem Open-Domain Datensatz, welcher Daten beinhaltet, welche allgemein gültig sind, und nicht nur in einer bestimmte Domäne.
Für jedes Dokument sind eine oder mehrere Sucheingaben definiert, mit dessen Eingabe das Dokument gefunden werden soll.
Darüber hinaus ist für jedes Dokument festgehalten, zu welchem Anwendungsfall sich dieses zuordnen lässt.
\myRef{Tabelle}{ausschnitt-studie-ergebnisse} zeigt einen Ausschnitt aus den Daten, welche durch die Studie generiert wurde.
Der Name des Projektes, aus welchem die Daten entnommen wurden, wurde unkenntlich gemacht.\\

\begin{table}[!ht]
    \centering
    \resizebox{\columnwidth}{!}{\begin{tabular}{|l|l|l|l|l|l|l|l|l|}
    \hline
        Anwendungsfall & Bereich & Sucheingabe & Erwartetes Dokument \\ \hline
        Onboarding & Entwicklung & Getting Started & Getting Started \\ \hline
    \end{tabular}}
\end{table}

\begin{table}[!ht]
    \centering
    \label{ausschnitt-studie-ergebnisse}
    \resizebox{\columnwidth}{!}{\begin{tabular}{|l|l|l|l|l|l|l|l|l|}
    \hline
        Gefundene Dokumente & Hit & Hit in Five & Hit in Three &  Hit in One \\ \hline
        getting started (legacy)..., onboarding usecase x..., onboarding re..., ui-debian-build-base..., ep... & true & true & true & true
    \end{tabular}}
    \caption[Ausschnitt der Ergebnisse]{Ausschnitt aus den Ergebnissen der Studie}
\end{table}

Die Suchfunktion gibt für jeden Algorithmus fünf Dokumente als Antwort auf eine Sucheingabe zurück.
Diese fünf Dokumente werden nach dessen Score sortiert.
Das bedeutet, dass das erste Dokument der Liste jenes ist, welches von dem Algorithmus als das passendste erachtet wird.
Die fünf Dokumente werden darauf untersucht, ob sich das gewünschte Dokument unter den Dokumente befindet.
Das Ergebnis ist ein Precision-Score für den Suchalgorithmus.
Um eine detailliertere Analyse zu ermöglichen wird nicht nur die Precision in Bezug auf die ersten fünf Dokumente gemessen.
Es wird die Precision für das erste Dokument (Hit in One), die ersten drei Dokumente (Hit in Three) und alle fünf Dokumente (Hit in Five / Hit) gemessen.
Anschließend werden die Precision-Scores der Algorithmen miteinander verglichen.
Dazu wird die Summe der Hits für alle Sucheingaben gebildet.
Die Summe wird durch die gesamte Anzahl der Sucheingaben dividiert.\\

Die Studie wird vollkommen automatisch durchgeführt.
Dadurch können Ergebnisse der Studie nicht durch Teilnehmer verfälscht werden.
Es bedeutet auch, dass die Studie eine hohe Reliabilität hat und mit jeder Durchführung das gleiche Ergebnis liefert.
Sie ist also reproduzierbar.
Der Nachteil dieser Herangehensweise ist, dass die subjektive Wahrnehmung des Nutzers, in Bezug auf die Precision der Suchfunktion, nicht beachtet werden kann.
So ist es denkbar, dass eine Sucheingabe nicht das gewünschte Dokument beinhaltet, aber andere Dokumente, welche ein Nutzer als sinnvoll erachten würde.
Die Ergebnisse der Studie sind damit abhängig von der Vorauswahl der Dokumente, welche gefunden werden sollen, und der Sucheingaben, welche a priori als sinnvoll bestimmt wurden.
Es ist möglich, dass eine Suchfunktion für eine Sucheingabe durchaus ein sinnvolles Ergebnis liefert, aber nicht das Ergebnis, welches durch den Aufbau der Studie erwartet wird.

\section{Auswertung der Ergebnisse}
Die Studie wurde mehrfach durchgeführt.
Die Precision-Scores werden auf zwei Nachkommastellen gerundet.
Die Precision-Scores werden in der \myRef{Tabelle}{precision-tabelle} zusammengefasst.

\begin{table}[!ht]
    \centering
    \label{precision-tabelle}
    \begin{tabular}{|l|l|l|l|}
    \hline
        Precision & Hit in Five & Hit in Three & Hit in One \\ \hline
        BM25 & 0,56 & 0,5 & 0,39 \\ \hline
        Hybrid & 0,56 & 0,52 & 0,35 \\ \hline
        Semantic & 0,15 & 0,13 & 0,07 \\ \hline
    \end{tabular}
    \caption[Precision]{Precision der Suchfunktionen}
\end{table}

Es ist auffällig, dass der Precision-Score der semantischen Suche wesentlich schlechter ist als der Precision-Score der BM25 Suche.
Dies könnte sich wie folgt erklären lassen.
Bei der Implementierung der Suche wurden alle H1-Header, H2-Header und Paragraphen eines Dokuments zusammengefasst in jeweils einen String für das entsprechende HTML-Tag.
Der Algorithmus der Suchfunktion vergleicht nun den String der Sucheingabe mit den Strings in den Dokumenten.
Also mit dem String des Titels, dem String der H1-Header, dem String der H2-Header und dem String der Paragraphen.
Weil die einzelnen Inhalte der HTML-Tags zusammengefasst werden, können diese Strings größer werden als die Sucheingabe, welche in der Regel zwei bis drei Wörter umfasst.
Dadurch, dass die Sätze einen großen Größenunterschied haben, fällt entsprechend auch der Score der Dokumente niedrig aus.
Da nun alle Scores niedrig sind kann die Suchfunktion nicht so leicht das passendste Dokumente finden.
Dadurch werden zum Teil weniger relevante Dokumente gefunden und der Precision-Score fällt besonders niedrig aus.\\

Auch das Postprocessing könnte eine Auswirkung auf den Precision-Score der semantischen Suche haben.
Für Bag of Words Modelle, wie BM25, ist es typisch ein Stemming vor der Indizierung durchzuführen.
Zum einen hat dies positive Auswirkungen auf den Precision-Score.
Zum sorgt dies dafür, dass verschiedene Tokens als das gleiche Wort betrachtet werden, auch wenn sie nicht in der gleichen morphologischen Form stehen.
Die Details zu dieser Problematik wurden bereits in Kapitel \ref{chap:volltext-indizierung} erläutert.
Für die semantische Suche könnte ein Stemming allerdings negative Auswirkungen haben.
Denn bei dem Training eines Sentence Transformers werden ganze Sätze betrachtet und nicht nur die Stammformen der einzelnen Wörter.
Das bedeutet, dass der Sentence Transformer die Wörter nicht mehr korrekt semantisch zuordnen kann.
Dies wirkt sich wiederum negativ auf den Precision-Score der semantischen Suche aus.\\

Die Hypothese ist also, dass der Precision-Score der semantischen Suche nur aus dem Grund gering ist, weil neben dem Title-Tag auch die weiteren Tags indiziert werden.
Um diese Hypothese zu überprüfen wurde die Studie ein weiteres Mal durchgeführt.
Dieses Mal wurden lediglich die Title-Tags indiziert.
Das Title-Tag ist für jedes Dokument nur einmal vorhanden.
Außerdem sind die Titelsätze der Dokumente wesentlich kürzer als ganze Paragraphen.
Mit dieser Konfiguration sollte der Größenunterschied des Sucheingabe-Strings und der Titel-Strings kein entscheidender Faktor mehr sein.
Das Ergebnis dieser Konfiguration ist in \myRef{Tabelle}{precision-tabelle-2} dargestellt.\\

\begin{table}[!ht]
    \centering
    \label{precision-tabelle-2}
    \begin{tabular}{|l|l|l|l|}
        \hline
        Precision & Hit in Five & Hit in Three & Hit in One \\ \hline
        BM25 & 0,54 & 0,5 & 0,39 \\ \hline
        Hybrid & 0,54 & 0,48 & 0,35 \\ \hline
        Semantic & 0,28 & 0,22 & 0,11 \\ \hline
    \end{tabular}
    \caption[Precision (2. Durchführung)]{Precision der Suchfunktionen (2. Durchführung der Studie)}
\end{table}

In dieser Konfiguration hat die semantische Suche eine bessere Precision.
Die Hit in Five Precision liegt bei 0,28 für die Konfiguration, welche nur den Titel der Dokumente indiziert, gegenüber 0,15 bei einer vollständigen Indizierung.
Die Verbesserung der Precision kann mit drei Faktoren begründet werden.
Zum einen die beiden Hypothesen, welche mit dieser Beobachtung untersucht werden sollten.
Also zum einen der Größenunterschied von Sucheingabe-String und den Strings, welche indiziert werden.
Zum anderen die Tatsache, dass sich das Postprocessing der Dokumente negativ auf die Fähigkeit des Sentence Transformers auswirkt, die Semantik der Wörter zu verstehen.
Außerdem besteht die Möglichkeit, dass die Inhalte der Dokumente leicht von den Titeln der Dokumente abweichen können.
Das würde bedeuten, dass für den gegebenen Versuchsaufbau der Titel besser geeignet ist, um die gewünschten Dokumente zu finden, als das gesamte Dokument.
Die Precision für die BM25-Suche ist bei vollständig indizierten Dokumenten besser als bei alleiniger Beachtung des Titels.
Lediglich die semantische Suche hat sich bei der zweiten Konfiguration verbessert.
Das bedeutet, dass sich der gezeigte Effekt der Verbesserung der Suchfunktion tatsächlich mit der Hypothese begründen lässt, dass die Länge der Strings ein zu berücksichtigender Faktor ist.
Da die semantische Suche weiterhin schlechter ist als die BM25-Suche, scheint es aber weitere Faktoren zu geben, welche zu berücksichtigen sind.\\

Eine mögliche Erklärung für die weiterhin schlechten Ergebnisse der semantischen Suche ist die Tatsache, dass es sich bei dem verwendeten Datensatz um eine Closed-Domain handelt.
Das bedeutet, dass der Textcorpus Fachbegriffe beinhaltet, welche der Transformer mit hoher Wahrscheinlichkeit nicht kennt.
Dementsprechend schwer fällt es dem Transformer folglich, passende Vektoren zu generieren, welche die Semantik von den domänenspezifischen Fachbegriffen abbilden.
In \myRef{Kapitel}{chap:ausblick} wird auf Ansätze in der Literatur eingegangen, welche dieses Problem versuchen zu lösen.

\section{Diskussion des Studienaufbaus}
Wie bereits beschrieben gibt der Recall an, wie viele der relevanten Dokumente gefunden wurden.
Die Precision gibt lediglich an, wie viele der Dokumente, welche gefunden wurden, relevant sind.
Eine optimale Suchfunktion würde per Definition alle Dokumente finden, welche relevant sind, und keine irrelevanten Dokumente.
Umgekehrt ist eine schlechte Suchfunktion eine Suchfunktion, welche keine relevanten Dokumente findet, sondern nur irrelevante.
Dabei spielt es für den Nutzer aber keine Rolle, ob irrelevante Dokumente gefunden wurden.
Die Tatsache, dass die relevanten Dokumente nicht gefunden wurden, machen die Suchfunktion für den Nutzer nicht benutzbar. 
Die Studie untersucht die Precision der Suchalgorithmen.
Auf Grundlage der obigen Argumentation zeigt sich, dass der Recall-Score wichtiger ist als der Precision-Score.
Denn ist der Recall gering, dann ist die Suchfunktion nicht benutzbar.
Eine Suchfunktion mit einem hohen Recall, aber eine niedrigen Precision könnte dagegen viele relevanten Dokumente finden, aber auch viele irrelevante Dokumente.
Solange die relevanten Dokumente in der Liste weiter oben dargestellt werden, ist diese Zusammensetzung aus Precision und Recall gut.\\

Die Studie hat die Precision für das erste, die ersten drei und alle fünf Dokumente gemessen.
Die Studie hat allerdings keinen Recall gemessen.
% TODO: Warum ist der Recall so wichtig?
Grund dafür ist die Tatsache, dass um den Recall messen zu können, alle relevanten Dokumente für eine Sucheingabe bekannt sein müssen.
Um für eine Studie a priori Sucheingaben zu bestimmen, und alle Dokumente, welche für diese Sucheingabe relevant sind, müsste der gesamte Datensatz bekannt sein.
Da der Datensatz mehrere hundert Dokumente beinhaltet ist dies nicht möglich.
Und auch wenn der gesamte Datensatz bekannt wäre, dann müsste trotzdem eine Entscheidung darüber getroffen werden, welche Dokumente relevant sind und welche nicht.
Das wäre wiederum eine subjektive Entscheidung, sodass die Validität des Ergebnisses anzweifelbar wäre.

% TODO: Welche Auswirkungen hatte die Technologie auf die Performance
\chapter{Erläuterung des Ergebnisses}

\chapter{Implementierung}

Die Architektur von Suchmaschinen besteht aus drei Komponenten.
Aus dem Crawling, der Indexierung und der Suchfunktion selbst.
Das Crawling ist zuständig für das Finden von Websites.
Die Indizierung ist zuständig für das optimale Speichern der Informationen der Websites, und die Suche ist zuständig für das Verstehen der Nutzeranfrage und die Abfrage der relevantesten Informationen aus dem Index, sowie dessen Verarbeitung und Darstellung.
Im Folgenden wird der Begriff Dokument verwendet, um die Dateien zu beschreiben, welche durch einen Crawler gesucht und durch den Index verarbeitet werden.
Unter Dokumenten können hierbei auch eine Website verstanden, welche durch einen Webcrawler durchsucht werden.

\section{Crawling}
Bei der Implementierung eines Systems für eine Suchfunktion benötigt das System zunächst einen Datensatz von Dokumenten, welche überhaupt grundsätzlich über die Suchfunktion gefunden werden können.
Dieser wird mithilfe eines Crawlers aufgebaut.
Ein Crawler ist ein Algorithmus, welcher ein Dokument als Startpunkt bekommt, und anhand dessen neue Dokumente findet.
Der Algorithmus analysiert dazu das Dokument auf Links, welchen der Algorithmus anschließend folgt.
Die neuen Dokumente werden durch den Index verarbeitet und wiederum auf neue Links analysiert.
Dieses Verfahren kann beliebig lange und beliebig rekursiv durchlaufen werden, um den Index zu erweitern.
Neben dem Crawling können Indizes befüllt werden, indem eine Liste von Dokumenten übergeben werden, welche dem Index hinzugefügt werden sollen.

\section{Indizierung}
Bei der Indexierung werden Wörter und Tokens mit Dokumenten assoziiert.
Dazu werden Wörter aus gecrawlten Dokumenten extrahiert und in die Datenbank geschrieben.
Den Dokumenten werden IDs zugeordnet und diese IDs werden den Wörtern zugeordnet.
Dem Index werden nun weitere Informationen hinzugefügt, wie die Wortfrequenz.
Die Wortfrequenz gibt an, wie oft ein Wort in einem Dokument vorkommt.
Es wird auch gespeichert, an welchen Stelle des Dokuments das Wort vorkommt, und auch in wie vielen Dokumenten ein Wort vorkommt.
Diese Art der Indizierung wird auch als invertierter Index bezeichnet, weil den Wörtern bzw.
Tokens die Dokumente zugeordnet werden, und nicht umgekehrt.
Wenn der Nutzer nun ein Keyword in die Suche eingibt, dann sind diese Keywords oft bereits im Index vorhanden, sodass die Dokumente, welche diese Keywords beinhalten einfach dem Index entnommen werden können.
Aufgabe der Suche ist es anschließend die Ergebnisse aufzubereiten.

\subsection*{Volltext-Indizierung}
Bei der Indizierung der Wörter besteht die Problematik, dass gleiche Wörter in unterschiedlichen Formen existieren können.
So stammen „Heizung“ und „heizen“ beide von dem gleichen Wortstamm „heiz“ ab.
Um bei der Indizierung Speicherplatz zu sparen, können Wörter auf diesen Wortstamm reduziert werden, damit sie als ein einziges Wort betrachtet werden können.
Die Bildung des Wortstamms wird auch als Stemming bezeichnet.
Beim Stemming kann es jedoch zu Overstemming und Understemming kommen.
Overstemming bedeutet, dass zwei Wörter, die eigentlich nichts miteinander zu tun haben, also nicht semantisch gleich sind, den gleichen Wortstamm besitzen und als ein Wort betrachtet werden.
Ein Beispiel hierfür sind die Wörter „Wand“ und „wandere“, wie in „ich wandere“.
Beide besitzen den Wortstamm „wand“ und werden entsprechend als ein Wort betrachtet.
Understemming bedeutet, dass zwei Wörter, die eigentlich etwas miteinander zu tun haben, also semantisch gleich sind, nicht den gleichen Wortstamm besitzen und dadurch als zwei verschiedene Wörter betrachtet werden.
Ein Beispiel hierfür sind die Wörter „absorbieren“ und „Absorption“, welche die Wortstämme „absorb“ und „absorp“ besitzen.
Es gibt Techniken zur Vermeidung solcher Probleme, wie der Einsatz vollständiger morphologischer Analysekomponenten.
Hierauf soll aber nicht weiter eingegangen werden.
Zur Implementierung eines Volltext-Index werden Dokumente und Wörter in einer m x n Matrix angeordnet, wobei m die Anzahl der Dokumente ist und n die Anzahl der Wörter.
Die Werte in dieser Matrix werden anhand einer ausgewählten Metrik bestimmt.
Eine oft verwendete Metrik ist das Verhältnis der Wortfrequenz mit der invertierten Dokumentfrequenz.
Die Wortfrequenz gibt dabei an, wie häufig das in dem Dokument vorkommt.
Die Dokumentfrequenz gibt an, in wie vielen Dokumente ein Wort vorkommt.
Das Verhältnis gibt damit an, wie charakteristisch das Wort für ein bestimmtes Dokument ist.
Wenn das Wort charakteristisch für ein Dokument ist, dann kommt es in diesem Dokument häufig vor, aber in anderen Dokumenten nur selten.
Wenn das Wort nicht charakteristisch für das Dokument ist, dann kommt es in anderen Dokumenten genauso häufig oder häufiger vor als in diesem Dokument.
Die Metrik wird auch als tf-idf-Wert bezeichnet.
TODO Rechnung besser verstehen
Wir können uns nun einen m-dimensionalen Raum vorstellen, in dem jedes Dokument eine Dimension darstellt.
Jedes Dokument liegt räumlich auf der Achse der eigenen Dimension.
Wenn der Nutzer nun 
Die Matrix lässt sich konzeptionell so verstehen, dass sie aus n Merkmalsvektoren besteht.
Verwenden wir die eben beschriebene Metrik, dann gibt jeder Merkmalsvektor eine Gewichtung an, wie charakteristisch ein Wort für die verschiedenen m Dokumente ist.
Nun kann der Nutzer Keywords in die Suche eingeben.
Diese Keywords 

\section{Suffix-Trees}
TODO

\section{N-Gramm}
Ein N-Gramm ist das Ergebnis der Spaltung eines Textes in N Fragmente.
Fragmente können dabei Buchstaben, Wörter, Phoneme, Morpheme etc. sein.
N-Gramme können bei der Indizierung verwendet werden und werden auch in Algorithmen von Stemming und Information Retrieval verwendet, außerdem werden durch N-Gramme Rechtschreibprüfungen oder die Suche nach ähnlichen Wörtern ermöglicht.
Besteht ein Text aus einem Fragment, dann wird von einem Monogramm geredet.
Besteht es aus zwei, wird es Bigramm genannt.\\

Um Algorithmen wie eine Rechtschreibprüfung zu ermöglichen, wird neben N-Grammen auch Ähnlichkeitsmaß benötigt, welches beschreibt, wie ähnlich sich zwei Terme sind.
Das Wort Term ist an dieser Stelle synonym mit Wort zu verstehen.
Der Dice-Algorithmus ist ein Ähnlichkeitsmaß für Terme.
Er zerlegt zwei Terme in dessen N-Gramme und bestimmt anhand dieser die Ähnlichkeit der beiden Terme.
Der Dice-Algorithmus lautet wie folgt:\\

TODO Equation\\

Nehmen wir beispielsweise die beiden Terme a=“Wort“ und b=“Wirt“.
Verwenden wir in diesem Beispiel Bigramme, dann ergeben sich für a folgende Bigramme: „\$W“, „Wo“, „or“, „rt“, „t\$“.
Für b ergeben sich: „\$W“, „Wi“, „ir“, „rt“, „t\$“.
Die Schnittmenge der beiden Mengen, also die Bigramme, welche sowohl in a als auch in b vorhanden sind, lauten „\$W“, „rt“, „t\$“.
Es sind also drei Stück.
Das bedeutet der Zähler der Funktion ergibt 2*3 also 6.
Die Anzahl der Bigramme in a und b sind jeweils 5, also ist der Nenner 5+5, also 10.
Dann ist das Ähnlichkeitsmaß der beiden Wörter bei der Verwendung von Bigrammen 6/10 (60\%).
Bei der Verwendung von Trigrammen besitzt a „\$\$W“, „\$Wo“, „Wor“, „ort“, „rt\$“, „t\$\$“ und b „\$\$W“, „\$Wi“, „Wir“, „irt“, „rt\$, „t\$\$“.
Es gibt drei gemeinsame Trigramme.
Das Ergebnis ist (2*3)/(6+6) also 50\%.
An diesem Beispiel lässt sich zeigen, dass sich das Ergebnis des Ähnlichkeitsmaß verändert, wenn ein anderes N-Gramm verwendet wird.

\section{Information Extraction}
Information Extraction ist der Prozess des Extrahierens von Informationen aus Dokumenten.
Dazu müssen zunächst die Dokumente gefunden werden, welche die passenden Informationen enthalten.
Dieser Schritt wird Information Retrieval genannt.
Das Extrahieren von Informationen mit einer Vielzahl von Algorithmen und erfolgt in mehreren Schritten.
Der erste Schritt von Information Extraction ist das Einteilen des Textes in Sätze.
Dieser Schritt wird Sentence Segmentation genannt.
Anschließend werden aus jedem Satz die einzelnen Wörter extrahiert (Word Tokenization) und kategorisiert als Nomen, Verb, Objekt, Pronomen etc. (Part-of-Speech Tagging).
Für das Part-of-Speech Tagging gibt es zwei Implementierungsansätze, rule-based anhand von Grammatikregeln und statistic-based.
Es folgt Syntactic Parsing , also das Generieren eines Syntaxbaumes, welcher den Text repräsentiert.
Zuletzt wird mithilfe von Coreference Resolution identifiziert, welche Stellen im Text sich auf die gleichen Entitäten beziehen.
Ein Beispiel für Coreference Resolution ist, dass ein System versteht, dass mit „Angela Merkel“ und „Merkel“ oder auch „A.
Merkel“ die gleiche Person gemeint ist.
Mithilfe dieser Schritte können Key Phrase Extraction, Named Entity Recognition, Entity Disambiguation, Relation Extraction und Event Extraction implementiert werden.
Key Phrase Extraction ist das Identifizieren von Schlüsselwörtern im Text, welche für diesen Text charakteristisch sind und welche den Text klassifizieren.
Ein Rezept für das Backen von Schokoladenkeksen könnte beispielsweise mithilfe von Schlüsselwörtern, wie „Backen“, „Kekse“, „Schokolade“ beschrieben werden.
Named Entity Recognition ist das Identifizieren von Personen, Orten und Datumsangaben.
Entity Disambiguation ist das Auseinanderhalten von Entitäten, welche die gleichen oder ähnliche Wörter besitzen.
Unter dem englischen Wort „apple“ lässt sich nicht nur ein tatsächlicher Apfel verstehen, sondern es könnte sich auch um das Unternehmen apple handeln.
Welche Bedeutung des Wortes die korrekte ist, erschließt sich aus dem Kontext.
Das ist die Aufgabe von Entity Disambiguation.\\

TODO Image\\

Teilbereiche von Informationsextraktion sind Named Entity Recognition, Named Entity Linking, Relation Extraction, Knowledge Base Reasoning.
Informationsextraktion ist für viele Anwendungsfälle der erste Schritt in einer größeren Pipeline.
Das könnte z.B. ein Chatbot sein oder ein Übersetzer.
Named Entity Recognition ist die Erkennung von Eigennamen.
Named Entity Linking ist der Prozess, eine bestimmte Entität eine einzigartige Identität zu geben.
Die Entität wird auf eine Target Knowledge Base gemappt, welche weitere Informationen über ebendiese Entität enthält.
Relation Extraction erkennt Beziehung zwischen Entities.
Die Beziehung kann durch ein Subjekt Prädikat Objekt Triple repräsentiert werden.
Abgrenzung von Information Extraction und Information Retrieval.
Information Retrieval selektiert relevante Dokumente.
Information Extraktion extrahiert Informationen aus einem gegebenen Dokument.
Die Klassifikation von Texten wird auch als Topic Modeling bezeichnet.
Es wird die Häufigkeit von Wörtern aus bestimmten Kategorien bestimmt und damit ein Thema zugeordnet.
Dieses Verfahren könnte zur Ermittlung von Labels verwendet werden.
Die Labels gehen dabei über die Inhalte des Textes hinaus.
Ein Algorithmus zum Topic Modeling ist Explicit Semantic Analysis (ESA).
Language Modelling ist ein Task von NLP, welcher es ermöglicht Texte zu übersetzen, Spellchecking, Spracherkennungen und Handschrifterkennung zu implementieren.\\

TODO Image

\section{Feature Location}

Um den Use-Case Feature Location zu implementieren kann ein gemeinsamer Index zwischen Code und Wissensdatenbanken verwendet werden.
Diese werden mithilfe von Word Embeddings in einer Vektordatenbank gespeichert.
Das hat den Vorteil, dass die Ähnlichkeit zwischen Dokumenten in der Wissensdatenbanken und Dateien, Klassen oder Methoden im Code bestimmt werden kann.
Wenn sich der Softwareentwickler also an einer bestimmten Stelle im Code befindet, und zu dieser Stelle weitere Informationen braucht, dann sorgt der gemeinsame Index dafür, dass Dokumente mit großer Ähnlichkeit in der Wissensdatenbank gefunden werden.
Der Softwareentwickler bekommt also die relevantesten Dokumente zu der Stelle Code, die er sich gerade ansieht. 

// TODO Quellen, Implementierungen finden (gibt es sowas schon?)
\chapter{Zusammenfassung und Ausblick}
In dem letzten Kapitel wurde der Vergleich zwischen der Confluence-Suche und der neuen Suchfunktion durchgeführt.
Die Ergebnisse sind ebenfalls diskutiert worden.
Dieses Kapitel die fasst die Arbeit zusammen und stellt einen Ausblick dar.

\section{Zusammenfassung}
Zweck der Arbeit war die Konzeption einer neuen Suchfunktion, welche nach festgelegten Kriterien besser ist als die Confluence-Suchfunktionen.
Dazu wurden in \myRef{Kapitel}{chap:anwendungsfaelle} Anwendungsfälle definiert, anhand dessen die Nutzung der Suchfunktionen konkretisiert wurde.
Dies hat den Rahmen geschaffen, um später einen Vergleich der beiden Suchfunktionen durchführen zu können.
Nach der Definition der Anwendungsfälle, wurde die Technologie der Suchfunktion von Confluence in \myRef{Kapitel}{chap:technologie-von-confluence} erläutert.
Damit wurde ein besseres Verständnis für die Suchfunktion von Confluence geschaffen.
Dieses bessere Verständnis war die Grundlage für die Konzeption einer neuen Suchfunktion in \myRef{Kapitel}{chap:konzeption-der-suchfunktion}.
Denn das Verständnis war notwendig, um Schwachstellen in der Confluence-Suche zu identifizieren und Verbesserungspotenziale zu verstehen.
In \myRef{Kapitel}{chap:konzeption-der-suchfunktion} wurde zuerst die Technologie der neuen Suchfunktion erläutert.
Während Confluence einen Volltext-Index verwendet, verwendet die neue Suchfunktion Sentence Transformer für die Indizierung.
Dazu wird eine Vektordatenbank benötigt, sowie der Scoring Algorithmus HNSW.
Auf Grundlage der Verbesserungspotenziale der Confluence-Suche und der dargestellten alternativen Technologien, wurde die neue Suchfunktion abgeleitet.
Dazu wurde argumentiert, warum sich die dargestellten Technologien als eine Alternative zu der bestehenden Confluence-Suche eignen sollen.
% TODO: Welche Quellen sind das:
Es wurde anhand verschiedener Quellen belegt, dass eine semantische Suche auf Basis von Sentence Transformers gute Ergebnisse für Open-Domain Datensätze erzielen. 
Dann wurde die Implementierung der neuen Suchfunktion in \myRef{Kapitel}{chap:implementierung} erklärt.\\

Nachdem die neue Suchfunktion implementiert wurde, wurden die Precision-Werte für einen angefertigten Datensatz zwischen den beiden Suchfunktionen verglichen.
Das nötige Wissen für diesen Vergleich wurde in \myRef{Kapitel}{chap:evaluationsmethoden} erläutert.
Es wurde ebenfalls auf alternative Möglichkeiten zur Messung der Performance von Suchfunktionen eingegangen.
Der angefertigte Datensatz leitete sich dabei aus den dargestellten Anwendungsfällen aus \myRef{Kapitel}{chap:anwendungsfaelle} ab.
Der Datensatz entstammte aus einer Wissensdatenbank eines Softwareprojektes.
Die Software soll im Gesundheitswesen eingesetzt werden und umfasste entsprechend Fachbegriffe aus dem Gesundheitswesen.
Darüber hinaus beinhaltet die Wissensdatenbank Definitionen, welche sich spezifisch auf das Softwareprojekt beziehen.\\

Die Durchführung der Studie wurde in \myRef{Kapitel}{chap:vergleich-der-suchfunktionen} erklärt.
Die Studie kam zu dem Ergebnis, dass die neue Suchfunktion eine wesentlich schlechtere Performance besitzt als die Confluence-Suche.
Begründet wurde dieses Ergebnis mit der Tatsache, dass es sich bei dem Datensatz, welcher in der Studie verwendet wurde, um eine Closed-Domain handelt.
Da der Sentence Transformer nicht fine-tuned wurde, kennt der Sentence Transformer die domänenspezifischen Fachbegriffe nicht und kann daher keine passenden Vektoren generieren, welche dessen Semantik wiederspiegeln.
Im Ausblick wird erläutert, welche Ansätze es in der Literatur gibt, um dieses Problem zu lösen.

\section{Ausblick}
\label{chap:ausblick}
Die Studie hat ergeben, dass die semantische Suche, so wie sie implementiert war, nicht besser ist als eine Suche auf Basis von BM25.
Der Kontext, in welchem die Suchen verwendet werden, ist eine Closed-Domain und beinhaltet domänenspezifische Fachbegriffe.
Probablistische Modelle unterscheiden sich in der Performance nicht zwischen Closed-Domain und Open-Domain.
Grund dafür ist, dass probablistische Modelle die Dokumente ohne ein vorheriges Training indizieren, anders als dies bei Sentence Transformern der Fall ist. 
So könnte der Schluss gezogen werden, dass sich probablistische Modelle besser für Closed-Domain Datensätze eignen.
Dennoch gibt es in der Literatur Ansätze, um die Probleme von Sentence Transformers mit Closed-Domain Datensätzen zu beheben.\\

Gururangan et. al. zeigen, dass eine weitere Phase des pre-trainings zur Adaption an eine Domäne Performance Verbesserungen mit sich bringt (\cite{Gururangan_Marasović_Swayamdipta_Lo_Beltagy_Downey_Smith_2020}).

Wang et. al. untersuchen TODO: (\cite{Wang_Thakur_Reimers_Gurevych_2022}.
Außerdem untersuchen Wang et. al. \textit{TSDAE: Using Transformer-based Sequential Denoising Auto-Encoderfor Unsupervised Sentence Embedding Learning}\cite{Wang_Reimers_Gurevych_2021}) von Wang et. al. untersuchen Möglichkeiten, um Sentence Transformer an eine Domäne anzupassen.
Dieser Vorgang wird als Domain Adaption bezeichnet.

% TODO: Ergänzen
% TODO: Zum einen Verbesserung durch Fine-Tuning möglich. Transformer auf Closed-Domain Daten trainieren?
% TODO: Sparse Vector Modelle ggf. besser für Closed-Domain? 


% ########################### BIBLIOGRAPHY ############################

{
\sloppy
\printbibliography[title={\label{chap:references}Literaturverzeichnis}]{} 
}

% ############################# APPENDIX ##############################

\part*{Anhang}
\newcounter{chapterAppendix}
\setcounter{chapterAppendix}{1}
\renewcommand{\thechapter}{\Alph{chapterAppendix}}
\chapter{Anhang}
\label{chap:appendix}

\section{docker-compose.yml File für Weaviate}
\begin{verbatim}
  version: '3.4'
  services:
    weaviate:
      image: semitechnologies/weaviate:1.18.3
      command:
        - --host
        - 0.0.0.0
        - --port
        - '2000'
        - --scheme
        - http
      ports:
        - "2000:2000"
      restart: on-failure:0
      environment:
        PROMETHEUS_MONITORING_ENABLED: 'true'
        QUERY_DEFAULTS_LIMIT: 20
        AUTHENTICATION_ANONYMOUS_ACCESS_ENABLED: 'true'
        PERSISTENCE_DATA_PATH: "/var/lib/weaviate"
        DEFAULT_VECTORIZER_MODULE: text2vec-transformers
        ENABLE_MODULES: text2vec-transformers,qna-transformers
        TRANSFORMERS_INFERENCE_API: http://t2v-transformers:8080
        QNA_INFERENCE_API: "http://qna-transformers:8080"
        CLUSTER_HOSTNAME: 'node1'
      volumes:
        - /var/weaviate:/var/lib/weaviate
    t2v-transformers:
      image: semitechnologies/transformers-inference:sentence-transformers-msmarco-distilroberta-base-v2
      environment:
        ENABLE_CUDA: 0
    qna-transformers:
      image: electra-qna
      environment:
        ENABLE_CUDA: 0  
\end{verbatim}

\section{Initialisieren des Schemas in Weaviate}
\begin{verbatim}
  WeaviateClass.builder()
    .className(DOCUMENT_CLASS)
    .properties(
      dataServiceHelper.buildProperties(
        mapOf(
          DOCUMENT_URL to WEAVIATE_TEXT_DATATYPE,
          TITLE_TAG to WEAVIATE_TEXT_DATATYPE,
          H1_TAG to WEAVIATE_TEXT_DATATYPE,
          H2_TAG to WEAVIATE_TEXT_DATATYPE,
          PARAGRAPH_TAG to WEAVIATE_TEXT_DATATYPE
        )
      )
    ).vectorIndexConfig(
      VectorIndexConfig.builder()
        .distance("l2-squared")
        .ef(100)
        .efConstruction(128)
        .build()
    )
    .invertedIndexConfig(
      InvertedIndexConfig.builder()
        .bm25(
          BM25Config.builder()
            .b(.5f)
            .k1(.5f)
            .build()
        )
        .build()
    )
    .vectorizer(VECTORIZER)
    .build()
\end{verbatim}

% ############################ DECLARATION ############################

% declaration of academic honesty / Eidesstattliche Erklärung
\chapter*{Eidesstattliche Erklärung}
\label{chap:declaration}

Hiermit versichere ich, dass ich die vorliegende Arbeit ohne Hilfe Dritter und nur mit den angegebenen Quellen und Hilfsmitteln angefertigt habe. Ich habe alle Stellen, die ich aus den Quellen wörtlich oder inhaltlich entnommen habe, als solche kenntlich gemacht. Diese Arbeit hat in gleicher oder ähnlicher Form noch keiner Prüfungsbehörde vorgelegen.
\vspace*{1cm}
\begin{flushright}\small \underline{\documentLocation, den \documentDate \hspace{0.4\textwidth}}\end{flushright}

\clearpage

% #####################################################################

\end{document}